% --- 题目 ---
\problem[P84-1 (24-2)]{已知函数 $\displaystyle f(x,y) = \begin{cases} (x^2+y^2)\sin\dfrac{1}{xy}, & xy \neq 0, \\ 0, & xy = 0, \end{cases}$ 则在点 $\displaystyle (0,0)$ 处 ( \quad )
\begin{tasks}(2)
  \task $\displaystyle \frac{\partial f(x,y)}{\partial x}$ 连续, $\displaystyle f(x,y)$ 可微.
  \task $\displaystyle \frac{\partial f(x,y)}{\partial x}$ 连续, $\displaystyle f(x,y)$ 不可微.
  \task $\displaystyle \frac{\partial f(x,y)}{\partial x}$ 不连续, $\displaystyle f(x,y)$ 可微.
  \task $\displaystyle \frac{\partial f(x,y)}{\partial x}$ 不连续, $\displaystyle f(x,y)$ 不可微.
\end{tasks}}
\ansat{290；【十年真题】 - 考点：多元函数微分学的基本概念 - 1}

% --- 题目 ---
\problem[P84-4 (20-2)]{关于函数
\[ f(x,y) = \begin{cases} \displaystyle xy, & xy \neq 0, \\ \displaystyle x, & y = 0, \\ \displaystyle y, & x = 0, \end{cases} \]
给出以下结论:
\begin{multicols}{2}
\begin{itemize}
    \setstretch{3}
    \item[\textcircled{1}] $\displaystyle \left. \frac{\partial f}{\partial x} \right|_{(0,0)} = 1$;
    \item[\textcircled{2}] $\displaystyle \left. \frac{\partial^2 f}{\partial x \partial y} \right|_{(0,0)} = 1$;
    \item[\textcircled{3}] $\displaystyle \lim_{(x,y) \to (0,0)} f(x,y) = 0$;
    \item[\textcircled{4}] $\displaystyle \lim_{y \to 0} \lim_{x \to 0} f(x,y) = 0$.
\end{itemize}
\end{multicols}
其中正确的个数为 ( \quad )
\begin{tasks}(4)
\task 4.
\task 3.
\task 2.
\task 1.
\end{tasks}}
\ansat{291；【十年真题】 - 考点：多元函数微分学的基本概念 - 4}


% --- 题目 ---
\problem[P87-4 (24-1)]{设函数 $\displaystyle f(u,v)$ 具有 2 阶连续偏导数, 且 $\displaystyle \mathrm{d}f(1,1) = 3\mathrm{d}u + 4\mathrm{d}v$.

令 $\displaystyle y = f(\cos x, 1+x^2)$, 则 $\displaystyle \left. \frac{\mathrm{d}^2 y}{\mathrm{d}x^2} \right|_{x=0} = \underline{\hspace{4em}}. $}
\ansat{291；【十年真题】 - 考点一：多元复合函数的偏导数与全微分的计算 - 4}


% --- 题目 ---
\problem[P87-5 (20-1)]{设函数 $\displaystyle f(x,y) = \int_0^{xy} \mathrm{e}^{xt^2} \mathrm{d}t$, 则 $\displaystyle \left. \frac{\partial^2 f}{\partial x \partial y} \right|_{(1,1)} = \underline{\hspace{4em}}. $}
\ansat{291；【十年真题】 - 考点一：多元复合函数的偏导数与全微分的计算 - 5}


% --- 题目 ---
\problem[P87-8 (19-2)]{设函数 $\displaystyle f(u)$ 可导, $\displaystyle z = y f\left(\frac{y^2}{x}\right)$, 则 $\displaystyle 2x \frac{\partial z}{\partial x} + y \frac{\partial z}{\partial y} = \underline{\hspace{4em}}. $}
\ansat{291；【十年真题】 - 考点一：多元复合函数的偏导数与全微分的计算 - 8}

% --- 题目 ---
\problem[P87-9 (19-3)]{设函数 $\displaystyle f(u,v)$ 具有 $\displaystyle 2$ 阶连续偏导数, 函数 $\displaystyle g(x,y) = xy - f(x+y, x-y)$, 求 $\displaystyle \frac{\partial^2 g}{\partial x^2} + \frac{\partial^2 g}{\partial x \partial y} + \frac{\partial^2 g}{\partial y^2}$.}
\ansat{291；【十年真题】 - 考点一：多元复合函数的偏导数与全微分的计算 - 9}

% --- 题目 ---
\problem[P87-10 (17-1,2)]{设函数 $\displaystyle f(u,v)$ 具有 2 阶连续偏导数, $\displaystyle y=f(\mathrm{e}^x, \cos x)$, 求 $\displaystyle \left. \frac{\mathrm{d}y}{\mathrm{d}x} \right|_{x=0}, \left. \frac{\mathrm{d}^2 y}{\mathrm{d}x^2} \right|_{x=0}$.}
\ansat{291；【十年真题】 - 考点一：多元复合函数的偏导数与全微分的计算 - 10}

% --- 题目 ---
\problem[P87-2 (25-3)]{已知函数 $\displaystyle z = z(x,y)$ 由 $\displaystyle z + \ln z - \int_y^x x \mathrm{e}^{-t^2} \mathrm{d}t = 1$ 确定, 则 $\displaystyle \left. \frac{\partial^2 z}{\partial x^2} \right|_{(1,1)} = \underline{\hspace{4em}}. $}
\ansat{292；【十年真题】 - 考点二：多元隐函数的偏导数与全微分的计算 - 2}

% --- 题目 ---
\problem[P87-3 (23-2)]{设函数 $\displaystyle z=z(x, y)$ 由方程 $\displaystyle \mathrm{e}^z + xz = 2x - y$ 确定, 则 $\displaystyle \left. \frac{\partial^2 z}{\partial x^2} \right|_{(1,1)} = \underline{\hspace{4em}}. $}
\ansat{292；【十年真题】 - 考点一：多元复合函数的偏导数与全微分的计算 - 3}

% --- 题目 ---
\problem[P89-变式2 (99-1)]{设 $\displaystyle y=y(x)$, $\displaystyle z=z(x)$ 是由方程 $\displaystyle z=xf(x+y)$ 和 $\displaystyle F(x,y,z)=0$ 所确定的函数, 其中 $\displaystyle f$ 和 $\displaystyle F$ 分别具有一阶连续导数和一阶连续偏导数, 求 $\displaystyle \frac{\mathrm{d}z}{\mathrm{d}x}$.}
\ansat{292；【方法探究】 - 考点二：多元隐函数的偏导数与全微分的计算 - 变式2}

% --- 题目 ---
\problem[P91-1 (22-1)]{设函数 $\displaystyle z=xyf\left(\frac{y}{x}\right)$, 其中 $\displaystyle f(u)$ 可导. 若 $\displaystyle x\frac{\partial z}{\partial x} + y\frac{\partial z}{\partial y} = y^2(\ln y - \ln x)$, 则 ( \quad )}
\begin{tasks}(2)
  \task $\displaystyle f(1)=\frac{1}{2}, f'(1)=0$.
  \task $\displaystyle f(1)=0, f'(1)=\frac{1}{2}$.
  \task $\displaystyle f(1)=\frac{1}{2}, f'(1)=1$.
  \task $\displaystyle f(1)=0, f'(1)=1$.
\end{tasks}
\ansat{294；【十年真题】 - 考点：已知偏导数问题 - 1}

% --- 题目 ---
\problem[P91-5 (24-2)]{设函数 $\displaystyle f(u,v)$ 具有 2 阶连续偏导数, 且函数 $$\displaystyle g(x,y) = f(2x+y, 3x-y)$$ 满足 $\displaystyle \frac{\partial^2 g}{\partial x^2} + \frac{\partial^2 g}{\partial x \partial y} - 6 \frac{\partial^2 g}{\partial y^2} = 1$.
\begin{enumerate}
\item[(1)] 求 $\displaystyle \frac{\partial^2 f}{\partial u \partial v}$;
\vspace{0.3em}
\item[(2)] 若 $\displaystyle \frac{\partial f(u,0)}{\partial u} = u \mathrm{e}^{-u}$, $\displaystyle f(0,v) = \frac{1}{50} v^2 - 1$, 求 $\displaystyle f(u,v)$ 的表达式.
\end{enumerate}}
\ansat{294；【十年真题】 - 考点：已知偏导数问题 - 5}

% --- 题目 ---
\problem[P93-1 (17-3)]{二元函数 $\displaystyle z = xy(3 - x - y)$ 的极值点是 ( \quad )
\begin{tasks}(4)
  \task $\displaystyle (0,0)$.
  \task $\displaystyle (0,3)$.
  \task $\displaystyle (3,0)$.
  \task $\displaystyle (1,1)$.
\end{tasks}}
\ansat{295；【十年真题】 - 考点一：多元函数的无条件极值 - 1}

% --- 题目 ---
\problem[P93-4 (23-1)]{求函数 $\displaystyle f(x,y) = (y - x^2)(y - x^3)$ 的极值.}
\ansat{295；【十年真题】 - 考点一：多元函数的无条件极值 - 4}

% --- 题目 ---
\problem[P93-6 (22-2)]{已知可微函数 $\displaystyle f(u,v)$ 满足
\[ \frac{\partial f(u,v)}{\partial u} - \frac{\partial f(u,v)}{\partial v} = 2(u-v)\mathrm{e}^{-(u+v)}, \]
且 $\displaystyle f(u,0) = u^2 \mathrm{e}^{-u}$.
\begin{enumerate}
\item[(1)] 记 $\displaystyle g(x,y) = f(x, y-x)$, 求 $\displaystyle \frac{\partial g(x,y)}{\partial x}$;
\item[(2)] 求 $\displaystyle f(u,v)$ 的表达式和极值.
\end{enumerate}}
\ansat{296；【十年真题】 - 考点一：多元函数的无条件极值 - 6}

% --- 题目 ---
\problem[P93-7 (22-3-仅数学三)]{设某产品的产量 $\displaystyle Q$ 由资本投入量 $\displaystyle x$ 和劳动投入量 $\displaystyle y$ 决定, 生产函数为 $\displaystyle Q=12x^{\frac{1}{2}} y^{\frac{1}{6}}$. 该产品的销售单价 $\displaystyle p$ 与 $\displaystyle Q$ 的关系为 $\displaystyle p=1 160-1.5Q$. 若单位资本投入和单位劳动投入的价格分别为 6 和 8, 求利润最大时的产量.}
\ansat{296；【十年真题】 - 考点一：多元函数的无条件极值 - 7}

% --- 题目 ---
\problem[P93-8 (21-3)]{求函数 $\displaystyle f(x,y) = 2\ln|x| + \frac{(x-1)^2 + y^2}{2x^2}$ 的极值.}
\ansat{296；【十年真题】 - 考点一：多元函数的无条件极值 - 8}

% --- 题目 ---
\problem[P93-9 (20-1,2,3)]{求函数 $\displaystyle f(x,y)=x^3+8y^3-xy$ 的极值.}
\ansat{296；【十年真题】 - 考点一：多元函数的无条件极值 - 9}

% --- 题目 ---
\problem[P93-2 (18-1,2,3)]{将长为 $\displaystyle 2\mathrm{m}$ 的铁丝分成三段, 依次围成圆、正方形与正三角形, 三个图形的面积之和是否存在最小值? 若存在, 求出最小值.}
\ansat{296；【十年真题】 - 考点二：多元函数的条件极值 - 2}

% --- 题目 ---
\problem[P93 (22-1)]{若当 $\displaystyle x \ge 0, y \ge 0$ 时, $\displaystyle x^2+y^2 \le k\mathrm{e}^{x+y}$ 恒成立, 则 $\displaystyle k$ 的取值范围是 \underline{\hspace{4em}}.}
\ansat{296；【十年真题】 - 考点三：多元函数在闭区域上的最值}

% --- 题目 ---
\problem[P95-变式 (04-1)]{设 $\displaystyle z=z(x,y)$ 是由
\[ x^2 - 6xy + 10y^2 - 2yz - z^2 + 18 = 0 \]
确定的函数, 求 $\displaystyle z=z(x,y)$ 的极值点和极值.}
\ansat{297；【方法探究】 - 考点三：多元函数在闭区域上的最值}